\documentclass[11pt,a4paper]{article}

\usepackage[utf8]{inputenc}
\usepackage[T1]{fontenc}
\usepackage[english]{babel}
\usepackage[a4paper,margin=2.1cm]{geometry}
\usepackage{microtype}
\usepackage{booktabs}
\usepackage{tabularx}
\usepackage{array}
\usepackage{xcolor}
\usepackage{enumitem}
\usepackage{listings}
\usepackage{hyperref}
\usepackage{fancyhdr}
\usepackage{titlesec}

\definecolor{ansblue}{HTML}{1E3A8A}
\definecolor{anslight}{HTML}{F3F6FC}
\definecolor{ansgray}{HTML}{5B6472}
\definecolor{anscode}{HTML}{F8FAFC}
\definecolor{ansgreen}{HTML}{0F766E}
\definecolor{ansred}{HTML}{B91C1C}

\hypersetup{
  colorlinks=true,
  linkcolor=ansblue,
  urlcolor=ansblue,
  pdftitle={ANS Part 1 Exam Prep for Hamza},
  pdfauthor={Yusuf Hamza Celebi},
  pdfsubject={CSE 341 Part 1 Exam Preparation},
  pdfkeywords={ANS, CSE341, programming languages, lexer, parser, EBNF, AST}
}

\setlength{\parindent}{0pt}
\setlength{\parskip}{5pt plus 1pt minus 1pt}
\setlist[itemize]{leftmargin=1.2em,itemsep=2pt,topsep=3pt}
\setlist[enumerate]{leftmargin=1.5em,itemsep=2pt,topsep=3pt}

\titleformat{\section}{\Large\bfseries\sffamily\color{ansblue}}{\thesection}{0.7em}{}
\titleformat{\subsection}{\large\bfseries\sffamily}{\thesubsection}{0.7em}{}
\titleformat{\subsubsection}{\normalsize\bfseries\sffamily}{\thesubsubsection}{0.7em}{}

\pagestyle{fancy}
\fancyhf{}
\fancyfoot[L]{\footnotesize\color{ansgray}CSE 341 ANS Part 1 Exam Prep \(\cdot\) Yusuf Hamza \c{C}elebi}
\fancyfoot[R]{\footnotesize\color{ansgray}\thepage}
\renewcommand{\headrulewidth}{0pt}

\newcommand{\code}[1]{\texttt{\detokenize{#1}}}
\newcommand{\important}[1]{\textbf{\color{ansred}#1}}
\newcommand{\design}{\textbf{Design intent}}
\newcommand{\actual}{\textbf{Actual D2 behavior}}
\newcommand{\studentName}{Yusuf Hamza \c{C}elebi}

\newcolumntype{Y}{>{\raggedright\arraybackslash}X}

\lstdefinelanguage{ANS}{
  morekeywords={let,function,return,for,in,if,elif,else,print,
    element,compound,reaction,reactants,products,formula,name,
    atomic_mass,valence,symbol,check,balanced,energy,exothermic,endothermic},
  sensitive=true,
  morecomment=[l]{//},
  morestring=[b]",
  morestring=[b]'
}

\lstdefinestyle{ans}{
  language=ANS,
  basicstyle=\ttfamily\small,
  keywordstyle=\bfseries\color{ansblue},
  commentstyle=\itshape\color{ansgray},
  stringstyle=\color{ansgreen},
  backgroundcolor=\color{anscode},
  showstringspaces=false,
  columns=fullflexible,
  keepspaces=true,
  breaklines=true,
  frame=single,
  rulecolor=\color{anslight},
  framerule=0.4pt,
  xleftmargin=4pt,
  xrightmargin=4pt
}

\lstdefinestyle{csmall}{
  language=C,
  basicstyle=\ttfamily\small,
  keywordstyle=\bfseries\color{ansblue},
  commentstyle=\itshape\color{ansgray},
  stringstyle=\color{ansgreen},
  backgroundcolor=\color{anscode},
  showstringspaces=false,
  columns=fullflexible,
  keepspaces=true,
  breaklines=true,
  frame=single,
  rulecolor=\color{anslight},
  framerule=0.4pt,
  xleftmargin=4pt,
  xrightmargin=4pt
}

\begin{document}

\begin{center}
  {\footnotesize\color{ansgray}Gebze Technical University \(\cdot\) CSE 341: Concepts of Programming Languages\par}
  \vspace{8pt}
  {\Huge\bfseries\sffamily ANS Part 1 Exam Prep\par}
  \vspace{4pt}
  {\Large Alchemist's Notation Script\par}
  \vspace{8pt}
  {\normalsize Student: \textbf{220104004926 \studentName}\par}
  {\normalsize Exam 1: \textbf{Thursday, May 14, 2026, 08:30} \(\cdot\) 90 minutes\par}
  \vspace{8pt}
  {\small Prepared from the submitted PDFs, D2 C code, D2/D3 tests, and Sebesta Chapters 1, 3, 4, and 5.}
\end{center}

\vspace{8pt}
\tableofcontents
\newpage

\section{How To Use This Document}

This document is written for Hamza as a crash-course defense guide for the ANS Part 1 exam. The instructor can ask both general programming-language questions and questions targeted to the components claimed in the contribution report.

\textbf{Main answer strategy:} always say whether you are talking about the ANS design, the EBNF grammar, the C implementation, or future Part 2 semantics. Part 1 is a lexer and parser with AST output. It does not execute, type-check, or chemically validate programs yet.

\subsection{Hamza's Claimed Areas}

The D6 contribution report gives Hamza these main ownership areas:

\begin{itemize}
  \item D1 4.1 Language Overview.
  \item D1 4.3 Syntax and EBNF.
  \item D2 build system, README, and command-line usage.
  \item D2 lexer and token definitions.
  \item D2 AST structures, AST dump, and cleanup.
  \item D3 valid example programs:
    \begin{itemize}
      \item \code{molar_mass_}\par\code{demo.ans}
      \item \code{combustion_}\par\code{analysis.ans}
      \item \code{multi_reaction_}\par\code{survey.ans}
    \end{itemize}
  \item Hamza D5 retrospective.
\end{itemize}

So expect questions about why ANS is a chemistry DSL, how the EBNF avoids ambiguity, how the lexer classifies tokens, how the AST is represented, how to build/run/test D2, and what the three valid programs demonstrate.

\section{ANS In One Minute}

\textbf{ANS} means \textbf{Alchemist's Notation Script}. It is a small domain-specific language for chemistry. It represents elements, compounds, reactions, formula literals, quantities, functions, lists, and control flow.

\begin{lstlisting}[style=ans]
element H { atomic_mass = 1.008, valence = 1, symbol = "H" };
compound water { formula = H2O, name = "Water" };

reaction combustion {
    reactants = CH4 + 2*O2,
    products  = CO2 + 2*H2O,
    check     = balanced,
    energy    = exothermic 890.3 kJ/mol
};
\end{lstlisting}

The source code looks close to chemistry notation. That is the point of the DSL: users write \code{H2O}, \code{CH4 + 2*O2}, and \code{molar_mass(water)} directly instead of hiding chemistry in strings or generic arrays.

\subsection{What Part 1 Actually Implements}

\begin{itemize}
  \item Reads an ANS source file.
  \item Converts source characters into tokens.
  \item Parses tokens with a hand-written recursive-descent parser.
  \item Builds an AST.
  \item Prints the AST with \code{--dump-ast}.
  \item Rejects malformed programs with line-numbered diagnostics.
\end{itemize}

\subsection{What Part 1 Does Not Implement}

\begin{itemize}
  \item No interpreter.
  \item No type checker.
  \item No real molar-mass calculation.
  \item No runtime output from \code{print}.
  \item No atom-balance checking yet.
  \item No semantic duplicate/missing-field validation yet.
\end{itemize}

\important{Exam-safe wording:} the design wants compile-time formula and balance validation, but D2 Part 1 only parses the syntax and creates an AST.

\section{Sebesta Chapter 1: Design Criteria}

ANS uses Sebesta's language-evaluation criteria directly. The core argument is that ANS chooses domain clarity over general-purpose flexibility.

\begin{tabularx}{\linewidth}{@{}p{3cm}Y Y@{}}
\toprule
\textbf{Criterion} & \textbf{Meaning} & \textbf{ANS answer} \\
\midrule
Readability & How easily code can be understood. & Chemistry words and notation are visible: \code{element}, \code{compound}, \code{reaction}, \code{H2O}. \\
Writability & How easily users can express solutions. & Domain tasks are short: declare compounds, write reactions, call \code{molar_mass}. \\
Reliability & How well the language prevents errors. & The design reserves places for formula, unit, type, and balance checks before execution. \\
Cost & Training, implementation, execution, maintenance, reliability cost. & Higher implementation cost because formulas, units, and reactions need custom parsing/semantics. \\
Generality & Usefulness outside the domain. & Sacrificed: ANS is intentionally chemistry-specific. \\
Orthogonality & Few primitives combined in regular ways. & Partly sacrificed: fixed domain declarations replace generic primitives. \\
\bottomrule
\end{tabularx}

\subsection{Strong Defense Answer}

\textbf{Question: Why did you choose a chemistry DSL?}

Because chemistry has notation that is meaningful on the page: formulas, reaction sides, coefficients, and units. A general-purpose language could calculate molar masses, but the source would not show the chemistry as clearly. ANS improves readability and writability for chemistry users, while accepting higher implementation cost and less generality.

\section{Syntax And EBNF}

Hamza's D1 ownership includes syntax and EBNF, so know this section well. EBNF is a formal way to describe valid source-program forms. In Sebesta terms, syntax is the form of expressions, statements, and program units; semantics is their meaning.

\subsection{Program Structure}

\begin{lstlisting}[style=ans]
element H { atomic_mass = 1.008, valence = 1, symbol = "H" };
compound water { formula = H2O, name = "Water" };
function molar_mass(compound c) { return 0.0; };
\end{lstlisting}

Top-level statements include element, compound, reaction, function, variable declaration, assignment, augmented assignment, print, expression statement, and if statement.

\subsection{Core EBNF Ideas}

\begin{tabularx}{\linewidth}{@{}p{3.5cm}Y@{}}
\toprule
\textbf{Notation} & \textbf{Meaning} \\
\midrule
\code{<name>} & Nonterminal, an abstract syntactic category. \\
\code{"if"} & Terminal, exact source text. \\
\code{{ X }} & Zero or more repetitions. \\
\code{[ X ]} & Optional item. \\
\code{X | Y} & Alternative choice. \\
\code{EOF} & End of file. \\
\bottomrule
\end{tabularx}

\subsection{Expression Precedence}

ANS encodes precedence with separate grammar levels:

\begin{tabularx}{\linewidth}{@{}p{4cm}p{4cm}Y@{}}
\toprule
\textbf{Level} & \textbf{Operators/forms} & \textbf{Example} \\
\midrule
Comparison & \code{== != >= <= > <} & \code{mass > 18.0} \\
Additive & \code{+ -} & \code{a + b - c} \\
Multiplicative & \code{* /} & \code{2.0 * 3.0} \\
Unary & \code{-} & \code{-3.14} \\
Postfix & call, index, member & \code{water.name}, \code{f(x)}, \code{a[0]} \\
Primary & literals, identifiers, lists, parentheses & \code{H2O}, \code{[H,O,C]}, \code{(x)} \\
\bottomrule
\end{tabularx}

Example answer: \code{1.0 + 2.0 * 3.0} parses as \code{1.0 + (2.0 * 3.0)} because multiplicative expressions are deeper than additive expressions.

\subsection{Reaction Grammar}

Reaction syntax is separate from normal arithmetic:

\begin{lstlisting}[style=ans]
reactants = CH4 + 2*O2
products  = CO2 + 2*H2O
\end{lstlisting}

Inside a reaction expression:

\begin{itemize}
  \item \code{+} separates species.
  \item \code{*} attaches a numeric coefficient to a species.
  \item Species may be formula literals like \code{CH4} or identifiers like \code{methane}.
  \item \code{2*O2} is required; \code{2O2} is not accepted.
\end{itemize}

\subsection{Ambiguity Handling}

\begin{tabularx}{\linewidth}{@{}p{4.2cm}Y@{}}
\toprule
\textbf{Issue} & \textbf{ANS solution} \\
\midrule
Arithmetic precedence & Separate expression levels. \\
Associativity & Binary operators are built left-to-right in parser loops. \\
Reaction vs arithmetic & Reactions use their own nonterminal. \\
Coefficient vs formula subscript & Require \code{2*O2}. \\
Formula vs identifier & Uppercase starts formula; lowercase/underscore starts identifier. \\
Keyword vs identifier & Hard keywords from lexer; contextual words handled by parser position. \\
Dangling else & All branch bodies are braced blocks. \\
Assignment target vs call & D1 excludes calls from assignable targets. \\
\bottomrule
\end{tabularx}

\section{Lexical Structure And Token Definitions}

Hamza's D2 ownership includes lexer and token definitions. The lexer is the front end of the parser. It reads source characters and returns one token at a time.

\subsection{Token Categories}

\begin{tabularx}{\linewidth}{@{}p{3.4cm}p{4cm}Y@{}}
\toprule
\textbf{Category} & \textbf{Examples} & \textbf{Rule} \\
\midrule
Identifier & \code{water}, \code{mass}, \code{compute_yield} & Starts lowercase or underscore. \\
Number & \code{20}, \code{1.008}, \code{6.02e23} & Integer, decimal, scientific notation. \\
String & \code{"Water"} & Double quoted. \\
Char & \code{'H'}, \code{'\n'} & Single quoted. \\
Formula & \code{H}, \code{O2}, \code{H2O}, \code{CH4} & Starts uppercase. \\
Unit literal & \code{g/mol}, \code{kg/mol}, \code{kJ/mol} & Checked before slash operator. \\
Hard keyword & \code{let}, \code{function}, \code{return}, \code{if} & Identifier-shaped text promoted by keyword table. \\
Operator & \code{==}, \code{!=}, \code{+=}, \code{+}, \code{/} & Longer operators before shorter ones. \\
Separator & \texttt{\{ \}}, \texttt{( )}, \texttt{[ ]}, \code{,}, \code{;}, \code{.} & Single-character tokens. \\
Comment & \code{// comment} & Skipped. \\
\bottomrule
\end{tabularx}

\subsection{Formula vs Identifier}

\begin{itemize}
  \item \code{water} is an identifier.
  \item \code{H2O} is a formula literal.
  \item \code{Water} is not a normal identifier in ANS because uppercase starts formula scanning.
\end{itemize}

This makes the language easier to parse and easier to read: program names and chemistry notation are visually separated.

\subsection{Unit Handling}

\begin{itemize}
  \item Slash units: \code{g/mol}, \code{kg/mol}, \code{kJ/mol} become \code{TOK_UNIT_LIT}.
  \item Single-word units: \code{mol}, \code{g}, \code{kg}, \code{kJ} are scanned as identifiers and interpreted as units only after a number.
\end{itemize}

\begin{lstlisting}[style=ans]
let amount = 5.0 mol;
let energy = 890.3 kJ/mol;
\end{lstlisting}

\textbf{Good exam phrase:} the lexer uses longest-match checks for slash units, and the parser handles single-word units by context.

\section{D2 Build System, CLI, And Architecture}

D2 is a C11 project. It has a Makefile and no external dependencies.

\subsection{Commands}

\begin{lstlisting}[style=ans]
make
make debug
make test
./ans <file.ans>
./ans --dump-ast <file.ans>
\end{lstlisting}

\begin{itemize}
  \item \code{make} builds the \code{ans} binary.
  \item \code{make debug} adds debug flags.
  \item \code{make test} runs valid programs expecting exit 0 and invalid programs expecting exit 1.
  \item \code{--dump-ast} parses and prints the AST.
\end{itemize}

\subsection{File Responsibilities}

\begin{tabularx}{\linewidth}{@{}p{4cm}Y@{}}
\toprule
\textbf{File} & \textbf{Role} \\
\midrule
\code{lexer.h} / \code{lexer.c} & Token types, token struct, lexer state, scanner implementation. \\
\code{parser.h} / \code{parser.c} & Parser state, recursive-descent parser, AST construction. \\
\code{ast.h} / \code{ast.c} & AST node kinds, node memory, child lists, AST printing, cleanup. \\
\code{error.h} / \code{error.c} & Shared diagnostic format and error flag. \\
\code{main.c} & CLI parsing, file reading, lexer/parser setup, AST dump, exit code. \\
\code{Makefile} & Build, debug, test, clean. \\
\bottomrule
\end{tabularx}

\section{AST Structures, Dump, And Cleanup}

Hamza's D2 ownership includes AST structures and AST dump. The AST is the parser's output tree.

\subsection{AST Node Shape}

\begin{lstlisting}[style=csmall]
typedef struct AstNode {
    NodeKind kind;
    int line;
    char *text;
    double num;
    struct AstNode **children;
    int n_children;
    int cap_children;
} AstNode;
\end{lstlisting}

Meaning:

\begin{itemize}
  \item \code{kind}: what the node is, such as \code{ND_ELEMENT_DECL}, \code{ND_BINARY}, \code{ND_CALL}.
  \item \code{line}: source line for debugging.
  \item \code{text}: name, operator, field key, formula, string, unit, etc.
  \item \code{num}: numeric payload for number nodes.
  \item \code{children}: dynamically grown array of child nodes.
\end{itemize}

\subsection{Memory Management}

\begin{itemize}
  \item \code{ast_new} uses \code{calloc}.
  \item \code{ast_set_text} copies source text into owned memory.
  \item \code{ast_add_child} grows the child array with \code{realloc}.
  \item \code{ast_free} recursively frees children, child array, text, and node.
\end{itemize}

\subsection{AST Dump Example}

For a reaction, the dump shows the grammar structure:

\begin{lstlisting}[style=ans]
ReactionDecl [combustion]
  ReactionField [reactants]
    ReactionExpr
      ReactionTerm
        Formula [CH4]
      ReactionTerm
        Number [2]
        Formula [O2]
\end{lstlisting}

This proves the parser recognized the syntax and structure. It does not prove the reaction is semantically balanced.

\section{Names, Binding, Scope, Lifetime}

Even though Enes owned D1 4.6, Hamza reviewed it and can still be asked about it.

\subsection{Identifiers}

Legal identifiers match:

\begin{center}
\code{[a-z_][A-Za-z0-9_]*}
\end{center}

ANS is case-sensitive. \code{water} is an identifier, but \code{H2O} is a formula literal.

\subsection{Static Scoping}

ANS uses static scoping. Names resolve by source-text nesting: current block, enclosing block/function, then top level.

\textbf{Why static?} Chemistry calculations should not depend on the call chain. Static scoping supports readability, reliability, and future compile-time semantic checking.

\textbf{Dynamic scoping problem:} a function could accidentally use a caller's local \code{water} or \code{formula}, making the same code mean different things depending on who called it.

\subsection{Lifetimes}

\begin{tabularx}{\linewidth}{@{}p{4.5cm}p{3.4cm}Y@{}}
\toprule
\textbf{Entity} & \textbf{Lifetime} & \textbf{Reason} \\
\midrule
Top-level elements, compounds, reactions & Static & Program chemistry model exists for whole run. \\
Top-level \code{let} variables & Static & Whole-program lifetime. \\
Function parameters & Stack-dynamic & Each call gets fresh values. \\
Local \code{let} variables & Stack-dynamic & Freed when block/call exits. \\
Lists and chemistry objects & Implicit heap-dynamic & Runtime-managed; no \code{new}/\code{delete}. \\
\bottomrule
\end{tabularx}

\section{D3 Valid Programs}

Hamza owned the three valid D3 programs, so know what each program demonstrates.

\begin{tabularx}{\linewidth}{@{}p{4.2cm}Y@{}}
\toprule
\textbf{Program} & \textbf{What it demonstrates} \\
\midrule
\code{molar_mass_}\par\code{demo.ans} & Element and compound declarations, typed function, loop over formula, \code{if/elif/else}, list literal, and function calls. \\
\code{combustion_}\par\code{analysis.ans} & Reaction declaration, formula literals, \code{check = balanced}, exothermic energy, typed reaction/compound parameters, nested calls, augmented assignment. \\
\code{multi_reaction_}\par\code{survey.ans} & Multiple reactions, formula-literal and identifier species, compound lists, loop counters, strings/chars, member access, repeated calls. \\
\bottomrule
\end{tabularx}

These programs parse successfully with D2. They do not execute in Part 1.

\section{Malformed Tests And Diagnostics}

Enes owned the malformed files, but the full project report includes them. Know the idea:

\begin{tabularx}{\linewidth}{@{}p{4.4cm}Y@{}}
\toprule
\textbf{Malformed case} & \textbf{Purpose} \\
\midrule
\code{missing semicolon} & Parser reports expected semicolon near the next token. \\
\code{unknown field} & Parser reports invalid field such as \code{protons}. \\
\code{illegal @} & Lexer reports unexpected character, then parser recovery reports syntax issue. \\
\code{unmatched brace} & Parser reaches EOF or premature token while expecting a closing brace. \\
\code{unterminated string} & Lexer reports the string error; parser then recovers through the broken construct. \\
\bottomrule
\end{tabularx}

The diagnostic format is:

\begin{center}
\code{<file>:<line>:<col>: error: <message>}
\end{center}

\section{Hamza D5 Retrospective Points}

Hamza's D5 says:

\begin{itemize}
  \item The lexer and parser seem to work and pass the current tests.
  \item More cases can appear later, so passing current tests is evidence, not proof.
  \item The hard part was keeping lexical structure consistent with EBNF and implementing it exactly.
  \item C was chosen because it is a main systems programming language and makes lexer/parser implementation explicit.
  \item For Part 2, the team plans to reconsider \code{print}, likely changing it from a special statement into a built-in function-style construct like \code{mol()}.
\end{itemize}

\section{Important Caveats}

\subsection{Semantic Checks Are Future Work}

\design: ANS wants compile-time formula validation and reaction balance checking.

\actual: D2 Part 1 only parses syntax. It accepts reaction syntax and builds an AST. It does not verify atom counts.

\subsection{Assignment Target Mismatch}

\design: D1 says function calls are expressions but not assignment targets, so \code{f(x) = 3.0;} should be invalid.

\actual: the current parser starts assignment parsing from \code{parse_postfix}, which includes calls. So D2 currently accepts \code{f(x) = 3.0;}. A good fix is to implement a separate \code{parse_assignable} that excludes call postfixes.

\subsection{AST Dump Is Evidence, Not Proof}

An AST dump confirms that the lexer, parser, and AST printer worked for that input. It does not prove every grammar path, semantic rule, or future interpreter behavior.

\section{Likely Exam Questions And Model Answers}

\subsection{Language Overview}

\textbf{Q: What is ANS?}

ANS is Alchemist's Notation Script, a chemistry DSL for declaring elements, compounds, reactions, and stoichiometric-style calculations. Part 1 implements the lexer, parser, AST dump, and syntax-error handling.

\textbf{Q: Why is it a DSL?}

Because chemistry notation is first-class syntax. Users write \code{H2O}, \code{CH4 + 2*O2}, \code{compound water}, and \code{check = balanced} directly. That is more readable for chemistry than generic strings or arrays.

\subsection{EBNF And Parsing}

\textbf{Q: How does your grammar encode precedence?}

It separates expression levels: comparison, additive, multiplicative, unary, postfix, primary. Lower levels bind tighter. Therefore \code{a + b * c} parses as \code{a + (b * c)}.

\textbf{Q: Why no left recursion?}

Recursive descent cannot handle left-recursive rules because the parser calls itself before consuming input. ANS uses repetition and parser loops instead.

\textbf{Q: How does ANS avoid dangling else?}

Every branch body is a braced block. The \code{else} has a clearly delimited preceding \code{if}/\code{elif} structure, so the classic unbraced ambiguity is avoided.

\subsection{Lexer And Tokens}

\textbf{Q: How does the lexer distinguish \code{water} and \code{H2O}?}

By first character. Lowercase or underscore starts an identifier. Uppercase starts a formula literal.

\textbf{Q: How is \code{kJ/mol} tokenized?}

The lexer probes compound unit literals before treating \code{/} as the division operator. So \code{kJ/mol} is one unit token, not \code{kJ}, \code{/}, \code{mol}.

\subsection{Build And AST}

\textbf{Q: How do you build and test D2?}

Run \code{make} to build, \code{./ans file.ans} to parse, \code{./ans --dump-ast file.ans} to print the AST, and \code{make test} to run the valid and invalid test suites.

\textbf{Q: What does \code{ast_free} do?}

It recursively frees child nodes, the children array, the text payload, and the node itself. This matters because the C implementation owns its memory directly.

\subsection{Scope And Future Semantics}

\textbf{Q: Static or dynamic scoping?}

Static scoping. Names resolve from the source text. That keeps chemistry calculations predictable and supports compile-time checking.

\textbf{Q: Does D2 check balanced reactions?}

No. D2 parses \code{check = balanced} and records it in the AST. Real atom-count validation is planned for the semantic/type-checking phase.

\section{Three Handwritten A4 Sheet Summary}

\subsection*{Sheet 1: ANS And Sebesta}

\begin{itemize}
  \item ANS = Alchemist's Notation Script, chemistry DSL.
  \item Domain constructs: \code{element}, \code{compound}, \code{reaction}, formulas, units.
  \item Part 1 = lexer/parser/AST/errors, not interpreter/type checker.
  \item Sebesta criteria: prioritize readability, writability, reliability.
  \item Sacrifice: cost, generality, some orthogonality.
  \item DSL defense: chemistry appears directly as \code{H2O}, \code{CH4 + 2*O2}.
\end{itemize}

\subsection*{Sheet 2: EBNF, Lexer, Parser}

\begin{itemize}
  \item EBNF: \code{[X]} optional, \code{{X}} repeat, \code{|} choice.
  \item Precedence: comparison < additive < multiplicative < unary < postfix < primary.
  \item Reaction grammar separate: \code{+} separates species; \code{*} attaches coefficient.
  \item Require \code{2*O2}, not \code{2O2}.
  \item Identifier starts lowercase/underscore; formula starts uppercase.
  \item Slash units before \code{/}; single units parsed after number.
  \item Recursive descent = one parser function per grammar level.
  \item Avoid left recursion.
\end{itemize}

\subsection*{Sheet 3: Code, AST, Tests, Caveats}

\begin{itemize}
  \item Commands: \code{make}, \code{make test}, \code{./ans}, \code{./ans --dump-ast}.
  \item Files: lexer, parser, AST, error, main, Makefile.
  \item AST node: kind, line, text, num, children, counts.
  \item AST memory: \code{calloc}, copy text, \code{realloc}, recursive free.
  \item Valid D3 programs: molar mass demo, combustion analysis, multi-reaction survey.
  \item Error format: file:line:col:error.
  \item Caveat: no execution/type checking/balance checking in Part 1.
  \item Caveat: parser currently accepts \code{f(x)=3}, although D1 says calls are not assignable.
\end{itemize}

\section{Final Defense Script}

When answering, use this structure:

\begin{enumerate}
  \item State the layer: language design, lexer, grammar, parser, AST, or future semantics.
  \item Give a concrete ANS example.
  \item Name the Sebesta concept.
  \item Say what D2 actually does.
  \item If there is a limitation, admit it precisely and propose the fix.
\end{enumerate}

Example:

\begin{quote}
The design says reaction balance should be checked before execution for reliability. In Part 1, D2 only parses \code{check = balanced} as a reaction field and stores it in the AST. The actual atom-count check belongs to the Part 2 semantic checker, where compound formulas and reaction sides can be resolved.
\end{quote}

\end{document}
